\documentclass[conference]{IEEEtran}
\IEEEoverridecommandlockouts
% The preceding line is only needed to identify funding in the first footnote. If that is unneeded, please comment it out.
\usepackage{cite}
\usepackage{amsmath,amssymb,amsfonts}
\usepackage{algorithmic}
\usepackage{graphicx}
\usepackage{textcomp}
\usepackage{xcolor}
\def\BibTeX{{\rm B\kern-.05em{\sc i\kern-.025em b}\kern-.08em
T\kern-.1667em\lower.7ex\hbox{E}\kern-.125emX}}
\begin{document}

    \title{Developing a Parallelized Genetic Algorithm Library using OpenMP, MPI, and CUDA\\
    }

    \author{
        \IEEEauthorblockN{Solomiia Gadiichuk}
        \IEEEauthorblockA{\textit{Faculty of Applied Sciences} \\
        Ukrainian Catholic University\\
        Lviv, Ukraine \\
        hadiichuk.pn@ucu.edu.ua}
        \and
        \IEEEauthorblockN{Bohdan Zasymovych}
        \IEEEauthorblockA{\textit{Faculty of Applied Sciences} \\
        Ukrainian Catholic University\\
        Lviv, Ukraine\\
        zasymovych.pn@ucu.edu.ua}
        \and
        \IEEEauthorblockN{Marta Kosiv}
        \IEEEauthorblockA{\textit{Faculty of Applied Sciences} \\
        Ukrainian Catholic University\\
        Lviv, Ukraine\\
        kosiv.pn@ucu.edu.ua}
        \and
        \IEEEauthorblockN{...}
        \IEEEauthorblockA{\textit{Faculty of Applied Sciences} \\
        Ukrainian Catholic University\\
        Lviv, Ukraine\\
        ...}
    }

    \maketitle

\begin{abstract}
        While Genetic Algorithms (GAs) are powerful optimization tools, their high computational cost limits scalability. This paper details the development of a genetic algorithm library that includes four parallelization strategies across shared-memory, distributed, and accelerator-based architectures to solve this problem. We describe a basic multi-threaded Generational GA and outline future library additions: a CUDA-accelerated GA, an MPI-based Island Model, an OpenMP Cellular GA, and a hybrid OpenMP/AVX Differential Evolution algorithm. The performance of this library will be tested using 3D mathematical surfaces and high-dimensional polygon image approximation tasks.
    \end{abstract}

    \begin{IEEEkeywords}
        Genetic Algorithms, Parallel Computing, Software Library, OpenMP, MPI, CUDA, Optimization
    \end{IEEEkeywords}

    \section{Introduction}
    Evolutionary algorithms are useful for solving complex optimization problems where traditional methods fail or take too much time. However, evaluating fitness functions for large populations over many generations requires a lot of computing power. 

    To solve this, different parallelization methods can be used, depending on the hardware and the specific algorithm. This project aims to develop a fast library that includes four different algorithms: a classic Generational GA, an Island Model GA, a Cellular GA, and Differential Evolution (DE). Using tools like OpenMP, OpenMPI, oneAPI/oneTBB, and CUDA, we will adapt these algorithms to the best hardware options (shared memory, distributed memory, and GPUs). 

    This paper presents the current progress of the library. It details a basic multi-threaded generational algorithm and outlines the theory and plans for future parallel models.

    \section{Background and Preliminaries}

    Genetic algorithms (GAs) are search and optimization methods inspired by natural evolution. Unlike standard algorithms that often get trapped in a local optimum, GAs use probability to explore large search spaces by evolving a population of possible solutions over several generations. This allows them to effectively escape local traps and find the true global maximum (or minimum) of a problem.

    \subsection{Representation and Fitness}
    Every possible solution in a GA is represented as a sequence of parameters called a \textit{chromosome} or \textit{genotype}. How this solution actually looks or works in the problem is the \textit{phenotype}. Each solution is evaluated using a \textit{fitness function}, $F(x)$, which gives a score based on how well it solves the problem. The main goal of the algorithm is to find a genotype with the highest (or lowest) fitness score.

    \subsection{Evolutionary Operators}
    Moving from one generation to the next uses three main steps:
    \begin{itemize}
        \item \textbf{Selection:} This step favors individuals with better fitness scores, giving them a higher chance to pass their genes to the next generation.
        \item \textbf{Crossover:} Selected pairs exchange parts of their chromosomes to create offspring. This combines the best traits from different parents.
        \item \textbf{Mutation:} To keep the population diverse and avoid getting stuck in local minimums, random changes are made to genes with a low probability.
    \end{itemize}

    \subsection{Schema Theory and Premature Convergence}
    GAs work by spreading short, successful gene patterns called \textit{schemas}. Over time, the algorithm combines these blocks to find the best solutions. However, a common problem in basic GAs is \textit{premature convergence}, where the population becomes too similar too fast and stops at a bad local solution. Methods like fitness scaling and different selection rules are used to control this and keep the population diverse.

    \section{Implemented Methodology}
    
    \subsection{Generational GA with Baseline Thread Parallelization}
    First, we implemented a basic Real-Coded Genetic Algorithm (RCGA). It uses continuous variables (double type) to find the minimum of mathematical functions. 

    The whole algorithm is written using C++ templates. This is very useful because it makes the library flexible and easy to reuse. We can easily apply different data types or new fitness functions without changing the main code.

    The algorithm uses the following main operations:
    \begin{itemize}
        \item \textbf{Initialization:} We create the first population randomly. We use local random number generators for each thread. This makes the code thread-safe and prevents race conditions.
        \item \textbf{Evaluation:} We calculate the fitness score for all individuals. This takes a lot of time, so we use OpenMP (\texttt{\#pragma omp parallel for}) to run it on multiple processor cores at the same time.
        \item \textbf{Selection:} We choose Tournament Selection. It is very fast because we do not need to sort the whole population. It is also easy to run in parallel.
        \item \textbf{Crossover:} We use Single-Point Crossover. It is simple and fast for the computer memory. It also keeps the exact good coordinates from the parents.
        \item \textbf{Mutation:} We use Gaussian Mutation. It adds a small random number to the genes. This helps the algorithm make small steps to find the exact minimum, and sometimes big jumps to escape bad local zones.
    \end{itemize}

    \section{Future Algorithmic Implementations}

    To build a complete parallelization toolkit, four advanced algorithms will be added to the library.

    \subsection{Generational GA with CUDA Acceleration}

    \subsection{Island Model GA with MPI}
        For cluster computers with distributed memory, the library will include an Island Model GA. Here, the main population is divided into isolated groups, or ``islands,'' that evolve on their own in separate processes. 

        To stop these isolated groups from getting stuck in local minimums, a migration step is added. Every $N$ generations, islands exchange some of their best individuals depending on the network setup (like a ring or fully connected graph). The library will use the Message Passing Interface (MPI) to handle this communication between computers.

    \subsection{Cellular GA}
        In this model, the population is organized as a two-dimensional grid where each cell represents an individual solution. Instead of interacting with the entire population, individuals interact only with their local neighbors defined by the Von Neumann neighborhood (up, down, left, and right). This spatial structure helps maintain diversity in the population and reduces the risk of premature convergence.
        
        Since the population is arranged as a grid, different cells can be processed independently. This structure makes the algorithm suitable for parallel processing where multiple cells of the population can be evaluated and evolved simultaneously.

    \subsection{Differential Evolution with Hybrid OpenMP/AVX}

    We plan to implement Differential Evolution (DE) for faster optimization of continuous variables. Unlike basic mutation, DE creates new solutions by calculating the mathematical difference between existing individuals in the population. This helps the algorithm adapt its search direction and converge much faster.

    \section{Benchmark Problems and Demonstrations}

    \subsection{3D Benchmark Surface (Swarm Convergence)}

    \subsection{Polygon Image Approximation}

    \section{Experimental Setup}
    % This section will detail the hardware specifications, cluster configurations, compiler versions, and specific GA parameters (population size, mutation rates, epoch lengths) used during testing.
    
    \section{Results and Discussion}
    % This section will present performance metrics, execution times, speedup graphs (Amdahl's Law analysis), and convergence rate comparisons to demonstrate the library's efficiency.

    \section{Conclusion}
    % This section will summarize the library's development progress, highlighting the capabilities of the implemented parallelization strategies across specific hardware constraints.

    \section*{Acknowledgment}

    \begin{thebibliography}{00}
        \bibitem{b1} % Add standard GA and Parallel computing references here later.
    \end{thebibliography}

\end{document}